\subsection{Measured hypothesis summary}
\label{sec:r-cross}
% Owner: A for eventual cross-network synthesis; B supplies the measured CIT cells.
\Cref{tab:cross} separates available CIT evidence from pending PS/CM experiments.
No ordering of network types by temporal sensitivity is established until the
same normalized measures have been computed for all applicable datasets.
The split H4 verdict avoids treating densification, diameter shrinkage,
descriptive fit and predictive superiority as one interchangeable claim.

\begin{table}[t]
\centering\small
\setlength{\tabcolsep}{4pt}
\begin{tabular}{@{}lcc p{3.8cm}@{}}
\toprule
ID & PS & CM & Measured CIT evidence \\
\midrule
H1 & Pending & Pending & Relative ranges: density/degree 1.798, clustering 2.361, LCC 1.114.\\
H2 & Pending & Pending & --\\
H3 & Pending & Pending & --\\
H4a & -- & -- & $\alpha=0.714$ [0.693, 0.730]; sublinear degree preference.\\
H4b & -- & -- & Lognormal preferred; scale 6.41 yr; $\Delta$AIC versus exponential 537.8.\\
H4c & -- & -- & $a=1.343$; $D_{90}=6$ throughout. Densification, no detected shrinkage.\\
H4d & -- & -- & Compact fits; MAE advantage over linear unresolved.\\
H5 & Pending & Pending & --\\
H6 & Pending & Pending & --\\
H7 & Pending & Pending & --\\
\bottomrule
\end{tabular}
\caption{Available evidence. Pending denotes unmeasured PS and CM experiments;
-- denotes inapplicability. CIT intervals are descriptive year resamples.
H4c/H4d have mixed verdicts. Values come from the shared summary and supporting tables.}
\label{tab:cross}
\end{table}
